\begin{statement}{E.75}{Proved}{The whole matrix source and the full Gaussian pro-three fibre product}
Retain the actual E.71 source and its finite image
\[
 J=K_6\cap\widehat F',\qquad
 \rho:J\twoheadrightarrow P_{71},\qquad
 |P_{71}|=2^{105}3^{17}.
 \tag{E.75.1}
Let $\psi:\widehat F'\twoheadrightarrow\mathcal H$ be the
six-evaluation Gaussian map of E.70. Thus $\mathcal H$ is the
specified class-two pullback of six copies of the free pro-three
group $P=\langle X,Y\rangle$, and
$\mathcal H^{\rm ab}\simeq\mathbf Z_3^{12}$.
Write
\[
 \mathcal G=(P/\Phi(P))^3\simeq\mathbf F_3^6,
 \qquad
 \epsilon:\mathcal H\longrightarrow\mathcal G
 \tag{E.75.2}
for projection to the first three factors followed by
abelianization modulo three. The other three translation values
are determined by the actual E.70 graph maps. The Gaussian
coordinate already retained by E.71 gives an epimorphism
$\pi:P_{71}\to\mathcal G$, and on the identical source $J$
one has $\pi\rho=\epsilon\psi$.

The entire common source image is exactly
\[
       (\rho,\psi)(J)=P_{71}\times_{\mathcal G}\mathcal H.
 \tag{E.75.3}
In particular the actual maximum common quotient is
$\mathcal G$, with the specified two maps in (E.75.2).
All six weight-three characters of $\mathcal H$ remain
independent in the kernel of the full E.71 projection.

Every one of the $55\,296$ indexed finite targets of E.74
has a compatible lift through the whole Gaussian pro-three
group and every one of its characteristic finite quotients,
retaining its complete E.71 value. The full Gaussian action
can be chosen to be the corresponding actual
$\operatorname{Inn}(x^j)$, $0\leq j<4$, with both peripheral
frames in $P$, as constructed in E.70.
For a chosen full correction value and one preimage $u\in J$,
the whole word fibre over the literal E.17 origin $f_i$ is
\[
       f_i u\bigl(\ker\rho\cap\ker(\psi|_J)\bigr).
 \tag{E.75.4}
The assertion retains the actual original word in each of the
twenty-four fibres.

There is also the following complete finite statement. Put
\[
 P_3=P/\overline{\gamma_4(P)\gamma_2(P)^3P^9},
 \qquad \mathbb G_3=P_3\rtimes\langle d\rangle,
 \quad |P_3|=3^7,
 \tag{E.75.5}
with the literal order-four action of E.70. Consider all
automorphisms $\alpha$ of $\mathbb G_3$ with frames
$h_y,h_z\in P_3$ satisfying
\[
 \alpha(x)=x,\qquad
 \alpha(y)=h_yyh_y^{-1},\qquad
 \alpha(z)=h_zzh_z^{-1}.
 \tag{E.75.6}
Require their six values
\[
 (h_y^{-1},h_y,h_z,h_z^{-1},h_z^{-1}h_y,h_y^{-1}h_z)
 \tag{E.75.7}
to belong to the actual image $S_3\subset P_3^6$ of
$\mathcal H$. There are exactly $324$ such automorphisms
and $972$ such framed targets. Exactly $81$ framed targets
lie above each of the twelve E.69 frames. Thus the entire
finite locus obtained by adjoining (E.75.5)--(E.75.7) to
E.71 has
\[
                 18\,432\cdot243=4\,478\,976
 \tag{E.75.8}
indexed values. This last count is finite. The whole-pro-three
automorphism assertion above uses the specified inner families
of E.70; it does not assign lifts of all $324$ finite
automorphisms.
\end{statement}

\begin{proof}
We identify the complete finite common image with the
twelve-character detector, prove that it detects every possible
three-group quotient of $P_{71}$, and then retain the full
Gaussian frames.

\paragraph{The exact whole-$J$ source.}
Adjoin the six Gaussian detector maps of E.70 to the entire
$11\,691$-point E.71 source. The new degree is $14\,283$.
The two lists of all $97$ literal $N_B$ words are compared
letter by letter. Compute their actual common image, and
then its entire derived subgroup, before making any
coordinate projection. The complete $N$ certificate has
$323$ strong generators and $8491$ word nodes, and its
independent replay checks all $228839$ final Schreier
relations.

Now retain the $7501$ E.71 coordinates and all $2592$
Gaussian coordinates. To the projected derived image
adjoin the $95$ actual sixth powers supplied by E.71's
integral lattice. Their original words are replayed in the
Gaussian maps and compared with the ordered products of
the original $97$ generator values, using the saved
unimodular columns. This retains exactly
\[
       J/\widehat N'=6\ker(e),
       \qquad e:\widehat N^{\rm ab}\longrightarrow
                         \widehat F^{\rm ab}.
 \tag{E.75.9}
In particular, a collection of word preimages that only
generates the preceding finite image is not substituted
for the complete new source.

The old E.71 complete base is placed first in the new
Schreier calculation. Its orbit product is checked to be
exactly $|P_{71}|$, and the remaining strong generators
give the actual kernel over that old image. Separate
complete certificates, each with an independent word
and relation replay, give
\[
 \begin{array}{c|c}
 \text{actual image}&\text{order}\\\hline
 S_3&3^{27}\\
 (\rho,\psi\bmod P_3^6)(J)&2^{105}3^{38}\\
 \text{kernel over }P_{71}&3^{21}.
 \end{array}
 \tag{E.75.10}
The finite common quotient therefore has order $3^6$.
The E.70 detector has the actual quotient map
\[
 \mathcal H/\Phi(\mathcal H)
       \xrightarrow{\ \sim\ }S_3/\Phi(S_3)
       \simeq\mathbf F_3^{12}.
 \tag{E.75.11}
In the pair order $(1,6),(2,5),(3,4)$, its four
coordinates per pair are the first entry's $X,Y$
abelian values and the two weight-three values of
the second entry times the inverse of the specified
lift of the first.

The new character certificate decodes every actual
six-factor matrix and every saved word in this quotient.
The kernel over $P_{71}$ has character image exactly
\[
 \{(0,0,a_1,a_2,\ 0,0,a_3,a_4,\ 0,0,a_5,a_6)
                      :a_i\in\mathbf F_3\}.
 \tag{E.75.12}
Its rank is six. The annihilator is precisely the span
of coordinates $0,1,4,5,8,9$ in zero-based numbering.
On every joint input and strong word, those coordinates
are compared with the six translation values read from
the actual E.71 Gaussian block. Thus (E.75.10) identifies
the particular common map $\pi$, rather than only its
group order.

For the source-word comparison, substitute the complete
derived and joint certificate DAGs into the original
$97$ generators and the actual $95$ sixth words. An
independent replay evaluates all $432$ joint basis
words, all $21$ old-kernel basis words, and the $95$
additional words on the entire $14\,283$ points. Its
$7539$ retained ancestors also give zero integral
$F$-exponents and zero $N$-exponents modulo six for
every one of these $548$ words. The kernel words kill
all $7501$ retained E.71 coordinates. Identity on
unretained coordinates is not needed for this map and
is not assigned.

\paragraph{Every finite three-group quotient of the old image.}
We use the literal target factors of E.71. On $J$, each
cyclic-cover deck exponent is zero. A retained matrix
block for $Q=\operatorname{GL}_2(\mathbf F_3)$ or
$\operatorname{SL}_2(\mathbf F_3)$ is exactly the right
action
\[
       (a,b):t\longmapsto a^{-1}tb,
       \qquad a=q\delta_n(w),\quad b=q(w).
 \tag{E.75.13}
The map from $Q\times Q$ to this permutation group
has kernel $\{(a,a):a\in Z(Q)\}$. Indeed, evaluation
at $t=1$ gives $a=b$, and evaluation at all $t$
then gives centrality. This kernel has order two.
Both groups $Q$ have cyclic Sylow-three subgroups of
order three, so the image of $Q\times Q$ has an
elementary abelian Sylow-three subgroup of order nine.
The coefficient block is the literal additive
$\mathbf Z/24$ action.

Each retained E.6 block is the actual direct product
$H_{96}\times C_6$, where $|H_{96}|=96$; its
Sylow-three subgroup is $C_3^2$. The E.62 factor is
a two-group. The E.67 factor retained by E.68 is
elementary abelian of order $3^9$, and E.71 verifies
that its entire $J$ projection remains this same
old image. The actual retained permutation map
therefore embeds $P_{71}$ in a product whose
Sylow-three subgroups are elementary abelian.
Every Sylow-three subgroup of $P_{71}$ is conjugate
into one of them, and is elementary abelian as well.

For any epimorphism from $P_{71}$ onto a finite
three-group, a Sylow-three subgroup maps onto that
quotient: its intersection with the normal kernel
is a Sylow-three subgroup of the kernel, and the
order formula gives surjectivity. Consequently
every finite three-group quotient of $P_{71}$ is
elementary abelian. No assertion about the entire
derived subgroup of $P_{71}$ is used.

\paragraph{The full pro-three common quotient.}
Put
\[
       T=\psi(J\cap\ker\rho)\ \triangleleft\ \mathcal H.
 \tag{E.75.14}
This is closed, since the source is compact. We have
$\psi(J)=\mathcal H$, because
$\widehat N'\subset J\subset\widehat F'$ and E.70
identifies both outer images with $\mathcal H$.
For $p\in P_{71}$, choose a source preimage $w\in J$
and send $p$ to $\psi(w)\bmod T$. This is a well-defined
continuous epimorphism
$P_{71}\to\mathcal H/T$. Thus $\mathcal H/T$ is a
finite three-group, and the preceding paragraph makes
it elementary abelian. It follows that
$\Phi(\mathcal H)\subset T$.

The image of $T$ in the finite detector is exactly
the kernel over $P_{71}$ in (E.75.10): both are
images of the same source subgroup $J\cap\ker\rho$.
Equations (E.75.11)--(E.75.12) therefore identify
$T/\Phi(\mathcal H)$ with the full kernel of the
ordinary six-coordinate map. Hence
\[
                       T=\ker\epsilon.
 \tag{E.75.15}
The actual commutative square $\pi\rho=\epsilon\psi$
then proves (E.75.3): choose an old preimage and
correct its Gaussian value by an element of the
displayed actual source kernel. This also proves
the specified maximum-common-quotient assertion.

Take now any indexed E.74 target. E.70 gives a
full Gaussian six-tuple $t\in\mathcal H$ reducing
to its finite frame, with action
$\operatorname{Inn}(x^j)$ and both full frames in
$P$. The required correction is
$\psi(f_i)^{-1}t$. Its $\epsilon$ value is the old
Gaussian coordinate already attached to the E.74
row. The fibre product (E.75.3) provides one actual
$J$ word with this full Gaussian value and the
specified entire E.71 value. This gives the claimed
compatible whole-tower lift. Equality of two such
word values is exactly the kernel condition in
(E.75.4).

\paragraph{The complete finite class-three frames.}
The calculation uses all $2187$ elements of $P_3$
in its exact Hall coordinates. For each $h_y$ put
\[
       y'=h_yyh_y^{-1},\qquad z'=(xy')^{-1}.
 \tag{E.75.16}
Retain precisely those choices with $(z')^4=1$.
The actual $P_3$-conjugacy class of $z$ has $729$
elements and each value has three frames $h_z$;
the calculation indexes this entire class and
retains every frame above every $z'$ in (E.75.16).

Here is the exact automorphism test. The original
discrete triangle is $C_2*C_4$, generated by $y,z$,
with the specified map to $C_4$. Its free kernel
has the literal basis $X=x^2y$, $Y=z^2y$. The
relations $(y')^2=(z')^4=1$ give an actual map of
that triangle, preserving the deck map. Completion
of its free kernel at three gives the map to
$P_3\rtimes C_4$. The verbal subgroup in (E.75.5)
is killed because its target satisfies those
relations. Thus this map factors through an actual
endomorphism of $\mathbb G_3$.
For every retained candidate, the two kernel images
\[
                         X'=x^2y',\qquad Y'=(z')^2y'
 \tag{E.75.17}
have determinant one on $P_3/\Phi(P_3)$. They
therefore generate $P_3$ by the finite three-group
Frattini criterion. The deck map is unchanged, so
the entire finite endomorphism is an automorphism.
Conversely every framed map (E.75.6) occurs in
this exhaustive list.

There are $324$ actions and $2916$ raw framed
choices. Apply the three actual class-two graph
conditions defining $S_3$ to (E.75.7); exactly
$972$ remain, with all $324$ actions still present.
For a Hall coordinate $X^aY^bc_0^e[c_0,X]^u[c_0,Y]^v$,
the exact reduction to E.69 is
\[
                     (a,b,e+2ab)\pmod3.
 \tag{E.75.18}
The discarded deck coordinate is checked to be
zero in every old record. This actual comparison
gives $81$ frames above each of the twelve E.69
frames. Each E.71 value fixes one of the four
old Gaussian translation targets, hence admits
$3\cdot81=243$ finite frames. The finite projection
of (E.75.3) realizes every one of them, proving
(E.75.8). The four specified actions
$\operatorname{Inn}(x^j)$, $0\leq j<4$, and their
twelve selected frames occur in this list. Their
whole-pro-three realization was proved in E.70;
no realization of all the other finite actions
is used here.

The new programs use the common prefix
\texttt{bc\_gaussian\_kn\_}, with suffixes
\texttt{common}, \texttt{characters}, \texttt{words}
and \texttt{frames}. The complete source, kernel,
word and character certificates retain every
projection map. The first frame-index run omitted
the explicit old deck-zero coordinate in comparing
record formats; that failed run and its exact
repair are preserved. A subsequent output-label
clarification specifies the range $0\leq j<4$
without claiming to enumerate all finite inner
automorphisms. No accepted E.70 or E.71 file was
changed.
\end{proof}
